\begin{explanationbox}
	\textbf{\textcolor{CExplanation}{一句话直觉}}
	仿射变换就是沿 $y$ 轴方向将椭圆“拉”成圆，把椭圆中的复杂关系转化为圆中的简单关系。

	\medskip
	\textbf{\textcolor{CExplanation}{核心拆解}}
	\begin{description}[style=nextline, leftmargin=2.8em, labelsep=0.5em, itemsep=0.45em]
		\item[\textbf{要点一（变换性质）：}]
			变换 $\varphi$ 是线性变换，只改变 $y$ 坐标：$x'=x,\ y'=\frac{a}{b}y$。它保持直线的平行性和分段比例，是将椭圆化为圆的线性拉伸变换。
		\item[\textbf{要点二（化圆过程）：}]
			将 $x=x',\ y=\frac{b}{a}y'$ 代入椭圆方程，恰好消去分母中的 $a,b$ 差异，得到圆方程 $x'^2+y'^2=a^2$。
		\item[\textbf{要点三（缩放因子）：}]
			$y$ 坐标拉伸倍数为 $\frac{a}{b}$，因此面积放大 $\frac{a}{b}$ 倍，斜率放大 $\frac{a}{b}$ 倍（因为纵坐标变化率放大）。
	\end{description}

	\medskip
	\textbf{\textcolor{CExplanation}{几何本质（线性拉伸）}}
	仿射变换 $\varphi$ 是坐标轴方向的缩放，不旋转，因此保持直线的平行性和共线点比例。它使椭圆变圆，同时保持图形的“线性结构”。

	\medskip
	\textbf{\textcolor{CExplanation}{代数意义（坐标变换）}}
	通过变量替换，将椭圆方程化为圆方程，使许多代数运算简化，积分类转化为圆中已知结论。

	\medskip
	\textbf{\textcolor{CExplanation}{考点价值（椭圆问题化圆）}}
	\begin{itemize}[leftmargin=*, itemsep=0.5em, topsep=0.4em]
		\item \textbf{考法一（面积转化）：} 椭圆内接三角形、四边形面积最值问题，可在圆中求解后按比例反求。
		\item \textbf{考法二（斜率计算）：} 椭圆中斜率乘积为定值、中点弦斜率等问题，可利用圆中垂径定理等转化。
	\end{itemize}

	\medskip
	\textbf{\textcolor{CExplanation}{顿悟点}}
	所有变换关系只依赖于一个比值 $a/b$，因此椭圆问题转化为圆问题，只需最后乘以或除以该因子即可。

	\medskip
	\textbf{\textcolor{CExplanation}{使用场景}}
	\begin{itemize}[leftmargin=*, itemsep=0.5em, topsep=0.4em]
		\item \textbf{场景一（面积最值）：} 求椭圆内接三角形最大面积时，变换为圆内接正三角形，面积乘以 $b/a$ 得椭圆面积。
		\item \textbf{场景二（斜率关系）：} 已知椭圆中两直线斜率之和或积，变换为圆中对应关系，利用圆的性质求解。
	\end{itemize}

\end{explanationbox}
